\section{Objects}

\begin{frame}[fragile]{Object Model Revision}
  C++ object lifecycle:
  \begin{itemize}
    \item Objects can be constructed in some storage location (their lifetime begins).
    \item When a destructor is called at that storage location, the lifetime ends.
  \end{itemize}

  For stack allocated objects:
  \begin{itemize}
    \item Storage location is known at compile time (offset within stack frame).
    \item Lifetime is known at compile time (when the scope ends).
    \begin{itemize}
      \item Lifetime begins when the constructor is called on the storage location.
      \item Compiler generates calls to the destructor at the end of the scope.
    \end{itemize}
  \end{itemize}
\end{frame}

\begin{frame}[fragile]{Objects and Methods}
  \begin{sqt:definition}
    A \emph{class} (in a general, language-agnostic sense) is approximately just a schema describing what data is present in an object.
    and a collection of functions that read/mutate objects satisfying that schema.
  \end{sqt:definition}

  \begin{sqt:definition}
    The C++ \emph{class} language feature lets you define the schema and the collection of functions ergonomically in one place.
  \end{sqt:definition}

  Methods (including constructors and destructors) implicitly take the \emph{address of the object} as an argument
  (binds to the argument \verb|this|, which is passed through the first argument register \verb|rdi|).
\end{frame}

\begin{frame}[fragile]{Construction}
  \begin{sqt:definition}
    Given a pointer to a chunk of memory, a \emph{constructor} is a function that
    initializes (starts lifetime of) an object at that location.\\
    It should be implemented such that the resulting object satisfies the \emph{class invariants}.
  \end{sqt:definition}

  Nit: lifetimes may start without any construction/initialization (namely implicit-lifetime types).
\end{frame}

\begin{frame}[fragile]{Analyzing Construction}
  \includesvg[inkscapelatex=false, width=\textwidth, height=0.8\textheight, keepaspectratio]{assets/constructor.drawio.svg}
\end{frame}

\begin{frame}[fragile]{New Expressions}
  \begin{sqt:definition}
    \emph{New expressions} are a language feature that allows for dynamic allocation and initialization of objects, in a simply and customizable way.
    A new expression looks up and calls the function \texttt{operator new()} (this is the customization point).
  \end{sqt:definition}

  A \emph{new expression} looks like \mintinline{cpp}{new (newArg1, newArg2) Gadget(ctorArg1, ctorArg2)}.
  When the compiler sees this, it performs a lookup for \mintinline{cpp}{operator new(size_t, ...)} which returns a pointer.
  It then calls the constructor of Gadget on the pointer.
\end{frame}

\begin{frame}[fragile]{New Expressions}
  \begin{minted}{cpp}
    /* 
    // automatically defined defaults, without having to include any headers!
    // these use malloc and free under the hood
    void* operator new(std::size_t) {...}
    ... operator new[], operator delete, etc
    */

    {
        // mallocs and constructs a Gadget on the heap
        auto x = new (newArg1, newArg2) Gadget(ctorArg1, ctorArg2);
    }

    {
        // mallocs space for a Gadget
        auto* tmp = static_cast<Gadget*>(operator new(sizeof(Gadget), newArg1, newArg2));
        // pseudocode: call Gadget constructor on the malloc result
        Gadget::Gadget(tmp, ctorArg1, ctorArg2);
        auto x = tmp;
    }
  \end{minted}
\end{frame}

\begin{frame}[fragile]{Operator New (Placement New Overload)}
  The language does not let you call \mintinline{cpp}{Gadget::Gadget(...)} on any 
  arbitrary chunk of memory and initialize an object there.
  
  \begin{sqt:definition}
    \emph{Placement new} is a workaround (uses header \code{<new>}) that allows
    you to call the constructor on any given address.
  \end{sqt:definition}

  \begin{minted}{cpp}
    new ((void *) 123456) Gadget{67};

    // gets transformed into approximately this pseudocode
    {
        // this operator new overload just returns 123456
        auto* tmp = static_cast<Gadget*>(operator new(sizeof(Gadget), 123456));
        // thus the constuctor is called on tmp
        Gadget::Gadget(tmp, 67);
    }
  \end{minted}
\end{frame}

\begin{frame}[fragile]{Destruction}
  \begin{sqt:definition}
    A destructor is a function that ends the lifetime of an object.
    It should be implemented such that it performs any necessary cleanup (e.g. freeing owned resources).
  \end{sqt:definition}

  \begin{sqt:result}
    For every initialization of an object that does not have \emph{dynamic storage duration} (i.e. not initialized via operator new),
    the compiler generates a call to this object's destructor at the appropriate location.\\
    All objects with dynamic storage duration must be destroyed manually.
  \end{sqt:result}
\end{frame}

\begin{frame}[fragile]{Destructor Generation Example}
  \begin{columns}[T]
    \begin{column}{0.45\textwidth}
      \begin{minted}[fontsize=\scriptsize]{cpp}
// what you write
void f(bool flag) {
    auto a = Gadget{1};
    {
        auto b = Gadget{2};
        auto c = Gadget{3};
    }
    if (flag) {
        auto d = Gadget{4};
        return;
    }
    auto e = Gadget{5};
}
      \end{minted}

      Destructor calls are generated at every scope exit, in \emph{reverse order of construction}.
    \end{column}
    \begin{column}{0.55\textwidth}
      \begin{minted}[fontsize=\scriptsize]{cpp}
// what the compiler generates (pseudocode)
void f(bool flag) {
    Gadget::Gadget(&a, 1);
    {
        Gadget::Gadget(&b, 2);
        Gadget::Gadget(&c, 3);
        Gadget::~Gadget(&c);
        Gadget::~Gadget(&b);
    }
    if (flag) {
        Gadget::Gadget(&d, 4);
        Gadget::~Gadget(&d);
        Gadget::~Gadget(&a);
        return;
    }
    Gadget::Gadget(&e, 5);
    Gadget::~Gadget(&e);
    Gadget::~Gadget(&a);
}
      \end{minted}
    \end{column}
  \end{columns}
\end{frame}

\begin{frame}[fragile]{Storage Re-use}
  Ending an object's lifetime does not release its storage.\\
  We are free to construct a different object in that storage.

  \begin{minted}{cpp}
    {
        auto g = Gadget{};  // lifetime of g begins
        g.~Gadget();        // lifetime of g ends, but its storage is still there

        auto* w = new (&g) Widget{}; // construct widget where g previously was
    } // g.~Gadget() called here (UNDEFINED BEHAVIOR)
  \end{minted}
\end{frame}

\begin{frame}[fragile]{Class Design Conventions}
  \begin{itemize}
    \item There are a bunch of \emph{conventions} to follow when implementing a class.
    \item Most programmers will assume these conventions hold when reading code.
    \item Most templated libraries (e.g. std::vector<T>) will assume that the type \code{T}
    you supply it obeys some of these conventions. Non-compliant types will lead to bugs.
  \end{itemize}

  A well designed class satisfies the following \emph{conventions}:
  \begin{itemize}
    \item Design by contract (class invariants, preconditions, postconditions)
    \item RAII and ownership semantics (acquire resources in ctor, release in dtor)
    \item Copy semantics (copy constructor creates an independent duplicate)
    \item Move semantics (move constructor transfers ownership of resources)
  \end{itemize}
\end{frame}

% TODO(jlu)
\begin{frame}[fragile]{Design by Contract, Well-Formedness}
  TODO(jlu) finish this slide
\end{frame}

% TODO(jlu) finish slide and come up with code example
\begin{frame}[fragile]{Ownership}
  Objects can \emph{own} state/resources or other objects.
  Objects can also simply have references to other objects.
  Always think about whether a relationship is an owning or non-owning relationship.

  TODO(jlu) finish this slide
\end{frame}

\begin{frame}[fragile]{RAII}
  \begin{sqt:definition}[Resource Acquisition is Initialization (RAII)]
    The name of a C++ programming technique/pattern/idiom which binds the life cycle of a resource to the lifetime of an object.
    It is done by acquiring the resource in the object's constructor, and releasing it in the destructor.
  \end{sqt:definition}

  Below is an example of an RAII class.

  \begin{columns}[T]
    \begin{column}{0.53\textwidth}
      \begin{minted}{cpp}
// Non-RAII design
struct IntBox {
    int *mData;
    void acquire(int x) { mData = new int(x); }
    void release() { delete mData; }
};
      \end{minted}
    \end{column}
    \begin{column}{0.47\textwidth}
      \begin{minted}{cpp}
// RAII design
struct IntBox {
    int *mData;
    IntBox(int x) : mData{new int(x)} {}
    ~IntBox() {delete mData;}
};
      \end{minted}
    \end{column}
  \end{columns}
\end{frame}

\begin{frame}[fragile]{Copy Construction/Assignment}
  \begin{sqt:definition}
    A copy constructor is a constructor that takes an existing object of the same class as its argument.
  \end{sqt:definition}
  Defining a copy constructor allows code like this to compile:
  \begin{minted}{cpp}
    /* 1.01 */
    auto a = Gadget{};
    auto& aLref = a;
    auto&& aRref = std::move(a);

    {
        // all these are copy constructions
        auto c = Gadget{a};
        auto d = Gadget{aLref};
        auto e = Gadget{aRref};
        // this is copy assignment, not copy construction
        e = a;
    }
  \end{minted}
\end{frame}

\begin{frame}[fragile]{Copy Construction}
  \begin{sqt:result}
    In a well-designed class, copy construction should create a "fully independent duplicate", such that
    both objects can be used and destroyed independently.
  \end{sqt:result}

  The compiler auto-generates a copy-ctor that recursively copies each member.
  \begin{minted}{cpp}
    struct BadIntBox {
      int *mData;
      BadIntBox(int x) : mData{new int(x)} {}
      ~BadIntBox() { delete mData; }
      // auto-generated copy constructor
      BadIntBox(const BadIntBox& other) : mData{other.mData} {}
    };
  \end{minted}
  Usually, automatically generated copy constructors are NOT well-designed.\\
  In this example, the resultant copy cannot be used nor destroyed independently.
\end{frame}

\begin{frame}[fragile]{Double Frees}
  Double frees occur when you write buggy classes that don't handle copying properly.

  Consider the below example using the \code{BadIntBox} we developed earlier
  \begin{minted}{cpp}
    int main() {
        auto original = BadIntBox{42};
        auto copy = BadIntBox{original}; // copy.mData = original.mData (both point to same heap object!)
    
        // ~original; ~copy; -> Double Free!
    }
  \end{minted}
  When the destructors run, both objects think that they are responsible for destroying
  the heap-allocated integer and freeing it. This results in Undefined Behavior!
  
  To fix this, think carefully when implementing your copy constructors.\\
  Make your class non-copyable if applicable.
\end{frame}

\begin{frame}[fragile]{Copy Construction}
  To perform a proper \emph{deep copy} on our \code{IntBox}, we want
  \begin{itemize}
    \item a new region of memory to be allocated on the heap
    \item the \code{other.mData} integer is copy constructed into this new region
  \end{itemize}

  \begin{minted}{cpp}
    struct IntBox {
      int *mData;
      IntBox(int x) : mData{new int(x)} {}
      ~IntBox() { delete mData; }
      IntBox(const IntBox& other) : mData{new int(*(other.mData))} {}
    };
  \end{minted}
\end{frame}
